\documentclass[11pt,a4paper]{article}
\usepackage{fontspec}
\setmainfont{DejaVu Sans}[Scale=0.90]
\usepackage{amsmath,amssymb}
\usepackage[margin=2.0cm]{geometry}
\usepackage{booktabs,array,tabularx}
\usepackage{graphicx}
\usepackage{enumitem}
\usepackage{titlesec}
\usepackage{parskip}
\titlelabel{\thetitle.\quad}
\titleformat{\section}{\normalfont\large\bfseries}{\thesection.}{0.6em}{}
\titleformat{\subsection}{\normalfont\normalsize\bfseries}{\thesubsection}{0.6em}{}
\titlespacing*{\section}{0pt}{14pt}{5pt}
\titlespacing*{\subsection}{0pt}{10pt}{4pt}
\setlist[itemize]{leftmargin=1.4em,itemsep=1pt,topsep=3pt}

\begin{document}

\begin{center}
{\Large\bfseries Identifiabilité sous les lois de conduction\\non-Fourier}\\[5pt]
{\normalsize Note de partie IV --- point de contrôle}\\[2pt]
{\small Version 2 --- 13 septembre 2026}
\end{center}

\vspace{4pt}

\section{Objet}

Krapez et Rigollet ont établi qu'en configuration revêtement sur substrat, sous excitation et
détection en face avant et en transfert unidimensionnel, l'épaisseur, la diffusivité et la
conductivité du revêtement ne sont pas identifiables séparément. La dégénérescence est qualifiée
d'\emph{inconditionnelle} : elle ne dépend ni du jeu de fréquences, ni du niveau de bruit, ni de la
méthode d'optimisation. Le cas des milieux gradués présente la même invariance, de dimension
fonctionnelle au lieu de un.

L'ensemble de ces travaux se place sous la loi de Fourier. La présente note examine ce que devient
cette dégénérescence sous les deux lois non-Fourier usuelles --- Cattaneo, puis Guyer--Krumhansl ---
et ce que les paramètres supplémentaires ainsi introduits apportent ou non.

\section{Cadre}

\subsection{Cattaneo}

La loi s'écrit $\tau_R\,\partial_t\varphi + \varphi = -\lambda\,\partial_z T$, d'où une conductivité
effective $\lambda_{\mathrm{eff}} = \lambda/(1+\tau_R p)$ après transformation temporelle.

\subsection{Guyer--Krumhansl}

\[ \tau_R\,\frac{\partial \mathbf{q}}{\partial t} + \mathbf{q}
= -\lambda\,\nabla T + \ell^2\left(\nabla^2\mathbf{q} + 2\nabla(\nabla\cdot\mathbf{q})\right) \]

En une dimension les deux termes non locaux se combinent en $3\,\partial_z^2 q$. L'élimination du
flux par le bilan d'énergie donne alors une conductivité effective dépendant de la fréquence,
\[ \lambda_{\mathrm{eff}}(p) = \frac{\lambda + 3\ell^2 p\,\rho c}{1 + \tau_R p} . \]

\subsection{Les deux combinaisons}

Dans les deux cas, la réponse de face avant ne dépend de la couche que par deux quantités :
\[ \sigma e = \xi_1 \sqrt{\frac{p\,(1+\tau_R p)}{1+\tau_\ell p}} ,
\qquad
\lambda_{\mathrm{eff}}\,\sigma = b \sqrt{\frac{p\,(1+\tau_\ell p)}{1+\tau_R p}} , \]
avec $\xi_1 = e/\sqrt{a}$, $b = \sqrt{\lambda\rho c}$, et
\[ \tau_\ell = \frac{3\ell^2}{a} , \]
temps de diffusion sur la longueur non locale. Cattaneo correspond à $\tau_\ell = 0$, Fourier à
$\tau_R = \tau_\ell = 0$.

\textbf{Les deux temps échangent leurs rôles} entre les deux combinaisons. Cette symétrie est à
l'origine du résultat de la section~6.

\subsection{Une substitution à écarter}

Remplacer $p$ par $p(1+\tau_R p)$ dans l'ensemble du quadripôle de Fourier donne le bon argument des
fonctions hyperboliques mais un coefficient de flux erroné d'un facteur $(1+\tau_R p)$. La phase
d'un milieu semi-infini part alors de $-45^\circ$ vers $-90^\circ$, tandis que la forme fermée de
Camacho de la Rosa \emph{et al.} donne $-45^\circ$ vers $0^\circ$.

La forme retenue ici reproduit cette dernière à la précision machine.

\section{Premier résultat : la dégénérescence n'est pas levée}

L'invariance établie sous Fourier est celle du groupe à un paramètre
\[ e \to \mu e, \qquad a \to \mu^2 a, \qquad \lambda \to \mu\lambda, \qquad \rho c \to \mu^{-1}\rho c , \]
sous lequel le temps de transit $\xi_1 = e/\sqrt{a}$ et l'effusivité $b$ sont invariants, donc le
quadripôle également. La forme infinitésimale en est l'identité d'Euler
\[ e\,\frac{\partial G}{\partial e} + 2a\,\frac{\partial G}{\partial a}
+ \lambda\,\frac{\partial G}{\partial \lambda} = 0
\qquad\text{pour toute fréquence.} \]

\textbf{Aucune longueur n'apparaît séparément dans la réponse.} Les quatre quantités dont elle
dépend sont $\xi_1$, $b$, et deux temps. Le temps de relaxation en est un ; la longueur non locale
n'entre que par $\tau_\ell = 3\ell^2/a$, qui en est un également. Or un changement d'échelle en
profondeur laisse les temps inchangés.

L'invariance survit donc aux deux lois. Sous Guyer--Krumhansl la relation d'Euler doit être étendue
d'un terme, la longueur non locale se transformant comme une longueur :
\[ e\,\partial_e G + 2a\,\partial_a G + \lambda\,\partial_\lambda G + \ell\,\partial_\ell G = 0 . \]

Cette transformation est cohérente avec la définition microscopique : $v^2 = 3a/\tau_R$ impose
$v \to \mu v$, donc $\ell = v\sqrt{\tau_N\tau_R/5} \to \mu\ell$.

La vérification numérique confirme la dérivation. Le résidu relatif de la combinaison d'Euler
demeure au niveau de l'erreur de troncature des différences finies, quelles que soient les valeurs
des deux temps. Le même calcul conduit dans le domaine temporel donne le même verdict, conformément
au principe de transposition énoncé par Krapez et Rigollet.

L'omission du terme en $\ell$ porte le résidu de $10^{-5}$ à $0{,}3$, ce qui établit que la longueur
non locale participe bien à l'invariance au lieu de la briser.

\textbf{Conclusion.} Ni un temps de relaxation fini, ni un terme non local ne restituent
l'identifiabilité perdue sous Fourier. La levée de l'indétermination continue d'exiger une
information imposant une échelle de longueur indépendante : mesure directe d'une propriété,
épaisseur mesurée par ailleurs, diffusion latérale, ou sources internes.

\section{Deuxième résultat : des paramètres authentiquement indépendants}

Les deux temps n'appartiennent pas à l'orbite d'invariance. Ils constituent donc des paramètres
supplémentaires dont l'estimation est une question distincte de celle des deux autres.

L'analyse de l'information de Fisher, conduite à épaisseur connue sur le triplet formé de la
conductivité, de la capacité volumique et du temps de relaxation, montre que l'introduction du
troisième paramètre laisse les incertitudes sur les deux premiers pratiquement inchangées. Il n'y a
donc pas de compétition entre eux : le temps de relaxation est porté par une direction de l'espace
des paramètres que les deux autres n'occupent pas.

\section{Troisième résultat : la loi d'échelle et son seuil}

La sensibilité de la phase mesurable au temps de relaxation se calcule en forme fermée. Pour un
milieu semi-infini en régime unidimensionnel,
\[ \varphi(\omega) = -\frac{\pi}{4} + \frac{1}{2}\arctan(\omega\tau) ,
\qquad
\frac{\partial\varphi}{\partial \ln\tau} = \frac{1}{2}\,\frac{\omega\tau}{1+(\omega\tau)^2} . \]

Cette sensibilité est maximale exactement en $\omega\tau = 1$, où elle vaut un quart de radian, et
décroît symétriquement de part et d'autre en échelle logarithmique.

L'incertitude relative d'estimation suit la loi correspondante. Le produit de l'incertitude relative
et de $\omega_{\max}\tau$ est constant sur plusieurs décades en deçà du seuil ; l'exposant local vaut
$-1{,}000$ à trois décimales, puis s'effondre au voisinage de l'unité. La constante dépend du niveau
de bruit et du nombre de points de mesure ; l'exposant et la position du seuil n'en dépendent pas.

\begin{center}
\includegraphics[width=0.52\textwidth]{02_omega_tau_threshold.png}
\end{center}

\textbf{Trois grandeurs distinctes ont leur point remarquable à la même abscisse.} L'argument de
l'énergie associée à la forme normale, $E = \tau\omega^2 - i\omega$, franchit la médiane de son
excursion à $-45^\circ$ en $\omega\tau = 1$. La phase mesurable franchit la sienne à $-22{,}5^\circ$
au même point. Et la sensibilité y est maximale, donc l'incertitude minimale.

La coïncidence n'est pas fortuite. La partie réelle du paramètre spectral, $\tau\omega^2$, est
précisément ce que la mesure doit percevoir pour distinguer la loi de Cattaneo de celle de Fourier ;
elle devient comparable à la partie imaginaire quand $\omega\tau$ atteint l'unité.

\section{Quatrième résultat : une cécité intrinsèque au matériau}

Les deux temps échangeant leurs rôles entre les deux combinaisons, leur égalité annule les racines :
\[ \tau_R = \tau_\ell
\qquad\Longrightarrow\qquad
\sigma e = \xi_1\sqrt{p} , \quad \lambda_{\mathrm{eff}}\sigma = b\sqrt{p} . \]

\textbf{La réponse est alors exactement celle d'un milieu de Fourier}, quelle que soit la valeur
commune des deux temps. La vérification numérique donne un écart de $3\times10^{-16}$, soit la
précision machine, sur quatre décades de temps de relaxation et huit de fréquence.

Un tel matériau est thermiquement indiscernable d'un matériau de Fourier. Aucune mesure thermique ne
peut le distinguer, et aucune amélioration du dispositif n'y change rien.

\subsection{Écriture microscopique}

Avec $\ell^2 = v^2\tau_N\tau_R/5$ et $a = v^2\tau_R/3$, il vient $\tau_\ell = \tfrac{9}{5}\tau_N$, de
sorte que la condition s'écrit
\[ \boxed{\;\tau_R = 1{,}8\;\tau_N\;} \]

C'est un rapport entre les temps de collision résistifs et normaux. Ces deux temps suivant des lois
de température différentes, un cristal peut traverser cette condition en refroidissant : il
paraîtrait de Fourier à une température précise, et non-Fourier de part et d'autre. Cette
conséquence est falsifiable.

\subsection{Critère de séparabilité}

\begin{center}
\includegraphics[width=0.66\textwidth]{04_guyer_krumhansl_blindness.png}
\end{center}

Au bruit de référence, un pour cent en amplitude et un dixième de degré en phase, les deux temps
doivent différer d'au moins $16\,\%$ pour une incertitude inférieure à $10\,\%$, et d'au moins
$25\,\%$ pour une incertitude inférieure à $1\,\%$. Loin de la diagonale l'incertitude est plate,
autour de $0{,}2$ à $0{,}4\,\%$ : il s'agit d'une crête étroite dans un plateau, non d'une
dégradation progressive.

\subsection{Comparaison avec la cécité d'interface}

Une seconde cécité a été établie par ailleurs : lorsque les effusivités du film et du substrat
coïncident, l'interface devient invisible et le temps de transit n'est plus mesurable.

Les deux conditions diffèrent par leur nature. Celle de l'interface dépend du choix de substrat,
donc elle se contourne. Celle-ci est intrinsèque au matériau et ne se contourne pas.

\section{Conséquence instrumentale}

Le seuil impose une condition sur la bande de mesure :
\[ \omega_{\max}\,\tau \;\gtrsim\; 1
\qquad\text{soit}\qquad
\tau \;\gtrsim\; \frac{1}{2\pi f_{\max}} . \]

Pour un banc de thermoréflectance en domaine fréquentiel dont la limite supérieure est de 200~MHz,
cette condition donne $\tau_{\min} = 8{,}0\times10^{-10}$~s. Cette valeur ne dépend que de
l'instrument, non du matériau.

Elle recoupe l'énoncé de Camacho de la Rosa \emph{et al.}, selon lequel les temps de relaxation
inférieurs à $3\times10^{-9}$~s se situent hors de la gamme de fréquences des montages courants.

\begin{center}
\begin{tabular}{@{}lccc@{}}
\toprule
Matériau & $\tau$ typique & $\omega_{\max}\tau$ à 200~MHz & Accessible \\
\midrule
Solides, borne basse & $10^{-14}$ s & $1{,}3\times10^{-5}$ & non \\
Solides, valeur courante & $10^{-12}$ s & $1{,}3\times10^{-3}$ & non \\
Solides, borne haute & $10^{-10}$ s & $0{,}13$ & marginalement \\
Gaz & $10^{-8}$ s & $12{,}6$ & oui \\
Milieux biologiques & $10^{-6}$ s & $1{,}3\times10^{3}$ & oui \\
\bottomrule
\end{tabular}
\end{center}

\textbf{Cas du nitrure d'aluminium.} Les temps de relaxation phononiques y sont de l'ordre de la
picoseconde à quelques dizaines de picosecondes, soit deux à trois ordres de grandeur en dessous du
plafond instrumental. Le temps de relaxation de l'AlN n'est pas accessible à une mesure de
thermoréflectance en domaine fréquentiel.

\section{Limites}

Les lois employées comportent un temps de relaxation unique par famille de processus, conformément à
l'approximation de Callaway. La dépendance en fréquence des temps de collision, qui est marquée dans
un cristal réel, est donc effacée : $\tau_R$ et $\tau_\ell$ sont des moyennes pondérées sur le
spectre.

La configuration traitée est celle d'une couche homogène sur substrat semi-infini. Le cas gradué
sous loi non-Fourier n'est pas examiné : la transformation de Liouville devrait y être reprise avec
une conductivité effective dépendant de la fréquence, et rien n'assure que la coordonnée transformée
reste indépendante de la fréquence. Le code refuse explicitement cette configuration.

La constante de la loi d'échelle dépend du niveau de bruit, du nombre de points et de leur
répartition. Seuls l'exposant et la position du seuil en sont indépendants.

Les vérifications numériques emploient le modèle qui a engendré les données. Elles établissent la
cohérence interne du calcul, non l'adéquation du modèle à une expérience réelle : ni la rugosité de
surface, ni la conductance d'interface, ni la phase de référence du montage n'y figurent, alors que
chacune produit un décalage de phase de la même nature que la signature recherchée.

Cette réserve n'est pas formelle. Une erreur de coefficient de flux a survécu à cent treize tests de
cohérence interne --- déterminant unimodulaire, limite homogène, composition, découpage --- parce
qu'elle était cohérente avec elle-même. Seule la confrontation à la forme fermée d'un article publié
l'a révélée. Toute loi implémentée exige désormais au moins une référence de validation extérieure.

\section{Références}

{\small
\begin{enumerate}[leftmargin=1.6em,itemsep=0pt]
\item J.-C. Krapez, F. Rigollet, \emph{Comment on ``Photothermal radiometry parametric identifiability
theory\ldots''}, arXiv:1708.07362 (2017).
\item J.-C. Krapez, \emph{Comment on ``Simultaneous density and thermal conductivity depth profile
reconstructions\ldots''}, J. Appl. Phys. \textbf{134}, 056101 (2023).
\item J.-C. Krapez, \emph{Linear, trigonometric and hyperbolic profiles of thermal effusivity in the
Liouville space and related quadrupoles}, Int. J. Therm. Sci. \textbf{136}, 182--199 (2018).
\item A. Camacho de la Rosa, R. Esquivel-Sirvent, D. Becerril, \emph{Relaxation times of non-Fourier
materials using frequency-domain thermoreflectance}, J. Appl. Phys. \textbf{137}, 155103 (2025).
\item D. Maillet \emph{et al.}, \emph{Thermal Quadrupoles}, Wiley (2000).
\end{enumerate}}

\vspace{4pt}
\noindent\small
Vérifications : \texttt{tests/test\_cattaneo.py}, \texttt{tests/test\_pulsed\_inversion.py}.
Figure : \texttt{notebooks/02\_omega\_tau\_threshold.py}. Seuils : \texttt{notes/conventions.md}.

\end{document}
